\documentclass{article}
\usepackage{xcolor}
\usepackage{stepdiff}

\stepdiffsetup{
  display=notes,
  theme=soft,
  layout=lecture,
  highlight=background,
  reason-style=muted,
  frame=true
}

\begin{document}

\[
\begin{stepdiff}
\step{x^2-1=(x-1)(x+1)}
\step[reason={expand}]{x^2-1=x^2-1}
\step[reason={final}, tag=final]{x^2-1=(x-1)(x+1)}
\end{stepdiff}
\]

\[
\begin{stepdiff}[display=focus, theme=minimal, layout=compact, highlight=underline, reason-style=plain, frame=false]
\step{a=b}
\step[reason={local options override setup}]{a=c}
\end{stepdiff}
\]

\stepdiffsetup{display=slide, theme=focus, highlight=color, reason-style=badge, show-reasons=false, frame=false}

\[
\begin{stepdiff}
\step{p=q}
\step[reason={hidden by setup}]{p=r}
\end{stepdiff}
\]

\[
\begin{stepdiff}[show-reasons=true, reason-style=muted]
\step{u=v}
\step[reason={shown locally}, tag=final]{u=w}
\end{stepdiff}
\]

\end{document}
